% ============================================================
\chapter{Measurable Functions}
\label{ch:measurable_functions}
% ============================================================

Before we can integrate, we must know which functions are allowed to be integrated. This is the question of \emph{measurability}: a function is measurable if the preimages of ``reasonable'' sets are themselves measurable. This condition, which may seem technical, is in fact very natural: it guarantees that one can unambiguously define the probability that $f(X)$ falls in a given interval, or the integral of $f$ with respect to a measure. Continuous functions are always measurable, but the class of measurable functions is vastly larger---and it is precisely this flexibility that makes the Lebesgue integral so powerful.

\section{Definitions and measurability criteria}

\begin{definition}[Measurable function]
Let $(X, \mathcal{A})$ be a measurable space. A function
$f : X \to \overline{\R} = [-\infty, +\infty]$ is \emph{$\mathcal{A}$-measurable}
(or simply \emph{measurable}) if:
\[
\forall\, a \in \R, \quad f^{-1}([-\infty, a)) = \{x \in X : f(x) < a\} \in
\mathcal{A}.
\]
\end{definition}

\begin{proposition}[Equivalent criteria]
\label{prop:measurability_criteria}
For $f : X \to \overline{\R}$, the following are equivalent:
\begin{enumerate}
    \item $\{f < a\} \in \mathcal{A}$ for all $a \in \R$;
    \item $\{f \leq a\} \in \mathcal{A}$ for all $a \in \R$;
    \item $\{f > a\} \in \mathcal{A}$ for all $a \in \R$;
    \item $\{f \geq a\} \in \mathcal{A}$ for all $a \in \R$.
\end{enumerate}
Moreover, if any of these holds, then $\{f = a\} \in \mathcal{A}$,
$\{f = +\infty\} \in \mathcal{A}$, and $\{f = -\infty\} \in \mathcal{A}$.
\end{proposition}

\begin{proof}
$(1) \Rightarrow (2)$: $\{f \leq a\} = \bigcap_{n=1}^{\infty} \{f < a + 1/n\}$.

$(2) \Rightarrow (3)$: $\{f > a\} = \{f \leq a\}^c$.

$(3) \Rightarrow (4)$: $\{f \geq a\} = \bigcap_{n=1}^{\infty} \{f > a - 1/n\}$.

$(4) \Rightarrow (1)$: $\{f < a\} = \{f \geq a\}^c$.

Finally, $\{f = a\} = \{f \leq a\} \cap \{f \geq a\}$,
$\{f = +\infty\} = \bigcap_n \{f > n\}$, $\{f = -\infty\} = \bigcap_n \{f < -n\}$.
\end{proof}

\begin{example}
\begin{enumerate}
    \item Every constant function is measurable.
    \item The indicator function $\mathbf{1}_A$ is measurable if and only if
    $A \in \mathcal{A}$.
    \item If $X$ is a topological space and $\mathcal{A} = \mathcal{B}(X)$, every
    continuous function $f : X \to \R$ is measurable.
    \item Every monotone function $f : \R \to \R$ is Borel measurable.
\end{enumerate}
\end{example}

\section{Operations on measurable functions}

\begin{theorem}[Algebraic stability]
\label{thm:algebraic_stability}
Let $f, g : X \to \R$ be measurable functions and $\alpha \in \R$. Then:
\begin{enumerate}
    \item $\alpha f$ is measurable;
    \item $f + g$ is measurable (when well-defined);
    \item $fg$ is measurable;
    \item $\abs{f}$ is measurable;
    \item $f/g$ is measurable (where $g \neq 0$);
    \item $f^+$, $f^-$, $\max(f, g)$, $\min(f, g)$ are measurable.
\end{enumerate}
\end{theorem}

\begin{proof}
(2) The key observation is:
\[
\{f + g < a\} = \bigcup_{r \in \Q} (\{f < r\} \cap \{g < a - r\}).
\]
This is a countable union of measurable sets, hence measurable.

(3) Use the identity $fg = \frac{1}{4}[(f+g)^2 - (f-g)^2]$ and the fact that $h^2$
is measurable when $h$ is (since $\{h^2 < a\} = \{-\sqrt{a} < h < \sqrt{a}\}$ for
$a > 0$).

(4) $\{\abs{f} < a\} = \{-a < f < a\} = \{f > -a\} \cap \{f < a\}$.

(6) $\max(f, g) = \frac{1}{2}(f + g + \abs{f - g})$ and
$\min(f, g) = \frac{1}{2}(f + g - \abs{f - g})$.
\end{proof}

\section{Limits of measurable functions}

\begin{theorem}[Stability under limits]
\label{thm:limit_stability}
Let $(f_n)_{n \in \N}$ be a sequence of measurable functions
$X \to \overline{\R}$. Then the following are measurable:
\begin{enumerate}
    \item $\sup_n f_n$ and $\inf_n f_n$;
    \item $\limsup_n f_n$ and $\liminf_n f_n$;
    \item if $f = \lim_n f_n$ exists (pointwise), then $f$ is measurable.
\end{enumerate}
\end{theorem}

\begin{proof}
(1) $\{\sup_n f_n > a\} = \bigcup_n \{f_n > a\} \in \mathcal{A}$. Similarly,
$\{\inf_n f_n < a\} = \bigcup_n \{f_n < a\} \in \mathcal{A}$.

(2) $\limsup_n f_n = \inf_n \sup_{k \geq n} f_k$ and
$\liminf_n f_n = \sup_n \inf_{k \geq n} f_k$, hence measurable by (1).

(3) If $f = \lim_n f_n$, then $f = \limsup_n f_n = \liminf_n f_n$.
\end{proof}

\begin{warning}[title=\textbf{Contrast with continuity}]
A pointwise limit of continuous functions need not be continuous (e.g.,
$f_n(x) = x^n$ on $[0,1]$). But a pointwise limit of measurable functions is
\emph{always} measurable. This is a major advantage of measurability over continuity.
\end{warning}

\section{Simple functions}

\begin{definition}[Simple function]
A function $\varphi : X \to \R$ is \emph{simple} if it takes only finitely many
values. It can be written as:
\[
\varphi = \sum_{k=1}^m c_k \mathbf{1}_{A_k},
\]
where $c_1, \ldots, c_m$ are the distinct values of $\varphi$ and
$A_k = \varphi^{-1}(\{c_k\})$ are the corresponding preimages. $\varphi$ is
measurable if and only if each $A_k \in \mathcal{A}$.
\end{definition}

\begin{theorem}[Approximation by simple functions]
\label{thm:simple_approximation}
Let $f : X \to [0, +\infty]$ be a measurable function. There exists a sequence
$(\varphi_n)_{n \geq 1}$ of measurable simple functions such that:
\begin{enumerate}
    \item $0 \leq \varphi_1 \leq \varphi_2 \leq \ldots$;
    \item $\varphi_n(x) \uparrow f(x)$ for every $x \in X$;
    \item if $f$ is bounded, the convergence is uniform.
\end{enumerate}
More generally, every measurable function $f : X \to \overline{\R}$ is the pointwise
limit of a sequence of measurable simple functions.
\end{theorem}

\begin{proof}
For $n \geq 1$ and $x \in X$, set:
\[
\varphi_n(x) = \begin{cases}
k \cdot 2^{-n} & \text{if } k \cdot 2^{-n} \leq f(x) < (k+1) \cdot 2^{-n},
\quad k = 0, 1, \ldots, n \cdot 2^n - 1, \\
n & \text{if } f(x) \geq n.
\end{cases}
\]
This function is simple (taking at most $n \cdot 2^n + 1$ values) and measurable
(each preimage is a measurable set of the form $\{k \cdot 2^{-n} \leq f <
(k+1) \cdot 2^{-n}\}$).

\emph{Monotonicity}: for every $x$, $\varphi_n(x) \leq \varphi_{n+1}(x)$ since the
partition at step $n+1$ refines that at step $n$.

\emph{Convergence}: if $f(x) < \infty$, for $n > f(x)$, we have
$0 \leq f(x) - \varphi_n(x) < 2^{-n} \to 0$. If $f(x) = +\infty$,
$\varphi_n(x) = n \to +\infty$.

\emph{Uniform convergence when $f$ is bounded}: if $f \leq M$, for $n > M$,
$\sup_x \abs{f(x) - \varphi_n(x)} \leq 2^{-n}$.

For general $f$, write $f = f^+ - f^-$ and approximate $f^+$ and $f^-$ separately.
\end{proof}

\begin{keyformulas}[title=\textbf{Approximation construction}]
\[
\varphi_n = \sum_{k=0}^{n2^n - 1} \frac{k}{2^n}\, \mathbf{1}_{\{k/2^n \leq f <
(k+1)/2^n\}} + n\, \mathbf{1}_{\{f \geq n\}}.
\]
\end{keyformulas}

\section{Measurable functions and Borel functions}

\begin{proposition}
Let $(X, \mathcal{A})$ be a measurable space. If $f : X \to \R$ is measurable and
$g : \R \to \R$ is Borel measurable, then $g \circ f$ is measurable.
\end{proposition}

\begin{proof}
For every $B \in \mathcal{B}(\R)$, $(g \circ f)^{-1}(B) = f^{-1}(g^{-1}(B))$.
Since $g$ is Borel, $g^{-1}(B) \in \mathcal{B}(\R)$. Since $f$ is
$(\mathcal{A}, \mathcal{B}(\R))$-measurable, $f^{-1}(g^{-1}(B)) \in \mathcal{A}$.
\end{proof}

\begin{corollary}
If $f$ is measurable, then $e^f$, $\sin(f)$, $f^p$ ($p > 0$), $\log(\abs{f})$
(where $f \neq 0$) are measurable.
\end{corollary}

\section{Almost everywhere convergence and convergence in measure}

\begin{definition}[Almost everywhere convergence]
We say $f_n \to f$ \emph{$\mu$-almost everywhere} ($\mu$-a.e.) if
\[
\mu\!\left(\{x \in X : f_n(x) \not\to f(x)\}\right) = 0.
\]
\end{definition}

\begin{definition}[Convergence in measure]
We say $f_n \to f$ \emph{in measure} if for every $\varepsilon > 0$:
\[
\mu\!\left(\{x \in X : \abs{f_n(x) - f(x)} > \varepsilon\}\right) \to 0
\quad (n \to \infty).
\]
\end{definition}

\begin{theorem}
If $\mu(X) < \infty$ and $f_n \to f$ a.e., then $f_n \to f$ in measure.
\end{theorem}

\begin{proof}
Set $A_n(\varepsilon) = \{x : \abs{f_n(x) - f(x)} > \varepsilon\}$ and
$B_N(\varepsilon) = \bigcup_{n \geq N} A_n(\varepsilon)$. Then
$B_N(\varepsilon) \downarrow \{x : \abs{f_n(x) - f(x)} > \varepsilon \text{ for
infinitely many } n\} \subset \{f_n \not\to f\}$.

By hypothesis, $\mu(\{f_n \not\to f\}) = 0$, so $\mu(B_N(\varepsilon)) \to 0$.
Since $A_N(\varepsilon) \subset B_N(\varepsilon)$,
$\mu(A_N(\varepsilon)) \to 0$.
\end{proof}

\begin{theorem}[Partial converse]
If $f_n \to f$ in measure, there exists a subsequence $(f_{n_k})$ such that
$f_{n_k} \to f$ a.e.
\end{theorem}

\begin{proof}
For each $k$, choose $n_k$ such that
$\mu(\{\abs{f_{n_k} - f} > 2^{-k}\}) < 2^{-k}$. By the Borel--Cantelli lemma
(Lemma~\ref{lem:borel_cantelli}), $\mu(\limsup_k \{\abs{f_{n_k} - f} > 2^{-k}\})
= 0$. Outside this set, $\abs{f_{n_k} - f} \leq 2^{-k}$ for $k$ large enough, so
$f_{n_k} \to f$.
\end{proof}

\begin{theorem}[Egorov]
\label{thm:egorov}
Let $\mu(X) < \infty$ and $f_n \to f$ $\mu$-a.e. For every $\delta > 0$, there
exists $A \in \mathcal{A}$ with $\mu(A^c) < \delta$ and $f_n \to f$ uniformly on $A$.
\end{theorem}

\begin{proof}
For $m, n \in \N^*$, set $E_{m,n} = \bigcup_{k \geq n} \{|f_k - f| \geq 1/m\}$.
For fixed $m$, $E_{m,n} \downarrow E_m$ where $E_m \subset \{f_n \not\to f\}$, so
$\mu(E_m) = 0$. Thus $\mu(E_{m,n}) \to 0$ as $n \to \infty$.

Choose $n_m$ so that $\mu(E_{m,n_m}) < \delta/2^m$. Set
$A = X \setminus \bigcup_m E_{m,n_m}$. Then
$\mu(A^c) \leq \sum_m \delta/2^m = \delta$. On $A$, for every $m$ and every
$k \geq n_m$, $|f_k - f| < 1/m$, which is uniform convergence.
\end{proof}

\section{Lusin's theorem}

\begin{theorem}[Lusin]
\label{thm:lusin}
Let $f : \R^n \to \R$ be Lebesgue-measurable and $\varepsilon > 0$. There exists a
closed set $F \subset \R^n$ with $\lambda^n(F^c) < \varepsilon$ and $f|_F$
continuous.
\end{theorem}

\begin{intuition}[title=\textbf{Interpretation of Lusin's theorem}]
``Every measurable function is nearly continuous.'' More precisely, one can make a
measurable function continuous by removing a set of arbitrarily small measure.
\end{intuition}

\section{Exercises}

\begin{exercise}
Show that if $f : \R \to \R$ is monotone, then $f$ is Borel measurable.
\end{exercise}

\begin{exercise}
Let $f : X \to \R$ be measurable and $g : X \to \R$ with $f = g$ $\mu$-a.e. Show
that $g$ need not be measurable (give a counterexample). However, show that if $\mu$
is complete, then $g$ is measurable.
\end{exercise}

\begin{exercise}
Show that convergence in measure is metrizable when $\mu(X) < \infty$ via the
distance $d(f, g) = \int_X \frac{\abs{f - g}}{1 + \abs{f - g}}\, d\mu$.
\end{exercise}

\begin{exercise}
Construct a sequence $(f_n)$ on $([0,1], \lambda)$ that converges in measure to $0$
but converges at no point. (\emph{Typewriter functions.})
\end{exercise}

\begin{exercise}
Let $(f_n)$ be a sequence of measurable non-negative functions. Show that
$\{x : \sum_n f_n(x) < \infty\}$ is a measurable set.
\end{exercise}
